\documentclass[12pt]{article}
\usepackage[utf8]{inputenc}
\usepackage[T1]{fontenc}
\usepackage{hyperref}
\title{Pumi, Northern Language Resources for K-12 Curriculum Development}
\author{Compiled by Claude Sonnet 5.5 (claude-sonnet-5-5) for Andr\'as Kornai}
\date{}
\begin{document}
\maketitle

\section{Introduction}
Northern Pumi (ISO 639-3 pmi; Glottocode nort2723) is a Sino-Tibetan language of China. Omniglot places it in Muli, Yanyuan and Jiulong counties in Sichuan and in Ninglang County in Yunnan. It has no standard orthography. The searches found no school materials for it. Most of what exists is scholarly description, and much of that covers other Pumi (Prinmi) varieties, which is noted below.

\section{Primary Educational Resources}
No primer, textbook or school programme for Northern Pumi was found. Omniglot reports that a pinyin-based Roman script has been proposed and that an attempt was made in Ninglang to teach children to write with the Tibetan script. No curriculum came from either.

\section{Linguistic and Reference Materials}
\begin{itemize}
\item Omniglot page on Pumi, with writing and dialect information: \url{https://omniglot.com/writing/pumi.htm}
\item Endangered Languages Catalogue, Northern Pumi bibliography: \url{https://endangeredlanguages.com/lang/5607/bibliography}
\item Glottolog record for Ding (2003), ``Prinmi'' in \emph{The Sino-Tibetan Languages}, tagged Northern Pumi: \url{https://glottolog.org/resource/reference/id/zurich:Ding2003}
\item Wikipedia article on the Pumi language (dialect groups after Ding 2014): \url{https://en.wikipedia.org/wiki/Pumi_language}
\item Ding Picus Sizhi, \emph{A Grammar of Prinmi} (Brill, 2014). It describes the Central dialect, not Northern Pumi, so it is a related reference only. It was seen in the Endangered Languages Catalogue listing above.
\end{itemize}
No dictionary specific to Northern Pumi was found.

\section{Religious and Community Resources}
An Asia Harvest people-group profile of Northern Pumi was found: \url{https://people-groups.asiaharvest.org/China/chinaPeoples/P/PumiNorthern.pdf}. It is a missionary-oriented demographic profile, useful only as background. No community language organisation was found.

\section{Digital and Audio Resources}
No audio, video, corpus or software for Northern Pumi was found.

\section{Curriculum Development Context}
The language is spoken by a large community but is not known to be taught in schools. A curriculum would need new work from the start: choosing an orthography, recording speakers, and writing materials. The scholarly descriptions above could guide phonology and grammar content.
\end{document}